We study whether large language models exploit an ``answer anyway'' loophole: when a request is impossible or the model is uncertain, does pressure to answer push the model toward compliant-looking but epistemically wrong output? Prior work has measured deception, hallucination, truthfulness, and abstention in separate benchmark traditions, but has not directly tested whether these behaviors share a common answer-pressure dynamic. We run a controlled cross-benchmark study on \simpleqa, \truthfulqa, and impossible false-presupposition prompts from \halogen. We evaluate two API models, \gptfourone and \gptfivemini, under three prompt conditions: a neutral baseline, an abstention-friendly instruction, and a forced-answer instruction. Across 960 model generations, the clearest effect appears on impossible prompts. When abstention is forbidden, both models attempt an answer on 100\% of impossible \halogen prompts, and most of those attempts are fabricated lists (96.7\% for \gptfourone and 91.7\% for \gptfivemini). On factual benchmarks, the effect is weaker and model-dependent. \gptfivemini shows a large abstention-versus-bluffing tradeoff: on \simpleqa, its attempt rate rises from 8\% to 100\% and its incorrect rate rises by 82 points under forced answering. \gptfourone, by contrast, tends to answer regardless of instruction. These results show that loophole-like behavior is real, especially when prompts reward surface compliance over epistemic honesty.
